\documentclass[12pt,a4paper]{article}

\usepackage[T1]{fontenc}
\usepackage[utf8]{inputenc}
\usepackage{lmodern}
\usepackage{microtype}
\usepackage{amsmath,amssymb}
\usepackage[margin=2.5cm]{geometry}
\usepackage{parskip}
\usepackage{titlesec}
\usepackage{enumitem}
\usepackage{xcolor}

\definecolor{slidegray}{RGB}{80,80,80}

\titleformat{\section}{\large\bfseries}{}{0em}{}[\vspace{-4pt}\rule{\linewidth}{0.4pt}]
\titlespacing{\section}{0pt}{16pt}{6pt}

\setlist[itemize]{itemsep=3pt,topsep=4pt}

\begin{document}

\begin{center}
  {\LARGE\bfseries Speaker Script}\\[6pt]
  {\large Can Simple SDEs Describe Equity-Market Time Series?}\\[4pt]
  {\color{slidegray} Dmitri Bolt\quad---\quad May 2026}
\end{center}

\bigskip

% ================================================================
\section{Slide 1. Title}
% ================================================================

This is an educational mini-research project.
My goal is not to build a production trading model.
I use textbook stochastic differential equation benchmarks from Øksendal
as a mathematical language and ask: which of these benchmarks describes
which transformation of real stock-price data, and in what sense?

% ================================================================
\section{Slide 2. Research Question and Financial Setup}
% ================================================================

I want to make sure the financial setup is clear, since we will be working
with real data.

A \textbf{stock} is a tradeable financial instrument --- a share of ownership
in a company.
Its \textbf{closing price} ($P_t^k$) is the last traded price at the end of each business day.
The market is open approximately \textbf{252 days per year} ---
weekends and holidays are excluded.

The key methodological point is that the same price series gives rise to five
structurally different empirical objects.
A model may be useful for one transformation and weak for another.
This is the organizing principle of the whole research.

% ================================================================
\section{Slide 3. Three Textbook Benchmarks}
% ================================================================

All three benchmarks come directly from Øksendal.

The \textbf{Itô diffusion} (Definition 4.1.1, equation 4.1.6) is the most general.
It says: the process changes by a deterministic drift term and a random diffusion term.
For daily log-returns, this is the natural language:
each day's return is a small noisy step.

The \textbf{geometric Brownian motion}
($\text{GBM}$, Chapter 5.1, equation 5.1.5)
specializes the drift to be proportional to the current level,
giving multiplicative growth.
In finance, N-sub-t ($N_t$) is the price level of a stock or a portfolio.
The solution is an exponential with a Brownian component ---
the standard picture of a random walk with drift for prices.

The \textbf{Ornstein--Uhlenbeck process}
($\text{OU}$, Exercise 5.7)
introduces mean reversion:
if the process drifts above its long-run level m,
the drift term ($\theta(m - X_t)$) pulls it back.
Here theta ($\theta$) is the mean-reversion speed.
This is exactly the behavior we expect from a detrended series
that oscillates around a long-run mean.

% ================================================================
\section{Slide 4. Data and Pipeline}
% ================================================================

We have 94 stocks, daily observations,
approximately 820 data points per stock.
The three representative stocks ---
NVDA (Nvidia), META (Meta Platforms), GOOG (Alphabet/Google) ---
are large US technology companies selected by explicit screening metrics
shown on the next slide.

The \textbf{detrending step} is the key operation:
we fit a simple ordinary least squares linear trend ($\text{OLS}$)
to each log-price series and keep the residual.
This residual ($\widetilde{X}_t^k$) ---
is the object we test against the Ornstein--Uhlenbeck benchmark.
Without detrending, stock prices drift upward over time
and the Ornstein--Uhlenbeck benchmark would clearly not apply.

% ================================================================
\section{Slide 5. Representative Stocks: Screening}
% ================================================================

I chose three stocks based on three criteria.

\textbf{Raw trend R-squared} ($R^2$):
how well does a linear trend explain the raw log-price?
A high value means the series is drift-dominated ---
good for the geometric Brownian motion benchmark.

\textbf{Lag-one autocorrelation function value}
($\text{ACF}$ at lag 1) of the detrended level:
how persistent is X-tilde ($\widetilde{X}_t^k$)?
A high value means slow mean reversion ---
good for the Ornstein--Uhlenbeck benchmark.

\textbf{Detrended standard deviation} ($\sigma$):
the amplitude of the detrended oscillations ---
larger means the visual picture is clearer.

NVDA has the strongest trend and largest amplitude.
META has the highest persistence of detrended levels.
GOOG is calmer but still well-expressed.
AAPL is weaker on all metrics and serves as a backup.

% ================================================================
\section{Slide 6. Raw Log-Prices vs Detrended Log-Prices}
% ================================================================

This is the core visual of the research.

The \textbf{left panels} show raw log-prices:
they grow over time, sometimes steeply,
and look exactly like geometric Brownian motion paths ---
drift-dominated with additive noise on top.

The \textbf{right panels} show the detrended residuals.
After removing the linear trend, the series oscillates around zero.
This is visually much more compatible with an Ornstein--Uhlenbeck benchmark.

The key insight: the same data yields both pictures depending on the transformation.

% ================================================================
\section{Slide 7. Increment Distributions: Log-Returns}
% ================================================================

The \textbf{log-return} ($\Delta X_{t_i}^k = \ln P_{t_i}^k - \ln P_{t_{i-1}}^k$) ---
is the standard measure of daily price change in finance.
It is symmetric and approximately scale-free:
a doubling of the price gives the same return regardless of the starting level.

The histograms show that detrending shifts the distribution center slightly toward zero
but does not reshape the noise.
The shape of the noise is the diffusion component;
the mean shift is the drift component.
The Itô diffusion language handles both.

% ================================================================
\section{Slide 8. Rolling Diagnostics: Is the Diffusion Constant?}
% ================================================================

The Itô benchmark (equation 4.1.6) with constant drift u and constant diffusion
coefficient v predicts that daily returns have constant mean and
constant standard deviation ($\sigma$) over time.

The 21-day rolling mean (left panel) stays small and fluctuates near zero ---
so the drift is hard to detect in daily data, which is consistent with noisy diffusion.

The 21-day rolling standard deviation (right panel) clearly changes over time:
\begin{itemize}
  \item NVDA: the ratio of maximum to minimum rolling volatility
        is approximately 3.7 times.
  \item META: the ratio is approximately 6.8 times.
  \item GOOG: the ratio is approximately 3.1 times.
\end{itemize}

This is a well-known empirical fact in finance called
\textbf{volatility clustering} ---
the market alternates between calm and turbulent regimes.
A constant-coefficient stochastic differential equation cannot capture this.
It remains a useful benchmark but has a known, quantifiable limitation.

% ================================================================
\section{Slide 9. OU Persistence: ACF of Detrended Levels}
% ================================================================

Ornstein--Uhlenbeck theory (Exercise 5.7) predicts that
the autocorrelation of the detrended series ($\widetilde{X}_t$) decays as ($e^{-\theta\ell}$),
where the mean-reversion speed is ($\theta$)
and the lag is ($\ell$).

The empirical autocorrelation function plots show exactly this pattern:
high positive lag-one values ---
0.9916 for NVDA, 0.9919 for META, 0.9880 for GOOG ---
that decay gradually rather than collapsing to zero immediately.

This does not prove an exact Ornstein--Uhlenbeck law.
But it is the empirical signature that motivates the Ornstein--Uhlenbeck approximation.
The slow geometric decay is consistent with theta-hat ($\hat\theta$) approximately 0.012,
corresponding to a half-life of about 57 trading days.

% ================================================================
\section{Slide 10. Detrended Levels and Rolling Mean}
% ================================================================

This slide shows the slow, undulating nature of the mean reversion.

The 63-day rolling mean (navy line) shows that the local long-run level of
X-tilde ($\widetilde{X}_t^k$) itself moves slowly over time.
This is consistent with very slow mean reversion:
the process does not snap back to a fixed equilibrium;
it drifts gently toward a slowly evolving local mean.

% ================================================================
\section{Slide 11. OU Parameter Estimation: AR(1) Proxy}
% ================================================================

The bridge between the continuous Ornstein--Uhlenbeck stochastic differential equation
and discrete daily data is the autoregressive model of order one (AR(1)) proxy.

If the process satisfies the Ornstein--Uhlenbeck equation, then the exact
discrete-time solution over one day is:
($\widetilde{X}_{t+1} = m(1-e^{-\theta}) + e^{-\theta}\widetilde{X}_t + \varepsilon_t$).

We estimate this as ($\hat\alpha + \hat\phi\,\widetilde{X}_t$).
From the fitted phi-hat ($\hat\phi$)
we recover theta-hat ($\hat\theta$) as minus the logarithm of phi-hat
($\hat\theta = -\ln\hat\phi$).

The \textbf{half-life} ($t_{1/2} = \ln 2\,/\,\hat\theta$)
is how long it takes for the deviation from the mean to halve.
For all three stocks the half-life is approximately 57 trading days ---
which is approximately \textbf{three calendar months}.
This is very slow mean reversion:
not a tight bounce, but a gentle multi-month drift back toward the long-run level.

% ================================================================
\section{Slide 12. Simulations: Connecting Theory and Data}
% ================================================================

This slide closes the loop between the theoretical benchmarks and the empirical data.

\textbf{Left:} pure Brownian motion --- symmetric random walk,
no drift, no mean reversion.

\textbf{Center:} geometric Brownian motion (GBM) ---
the same Brownian noise multiplied by the current level,
giving exponential growth under random perturbations.
This is the visual model for raw stock prices.

\textbf{Right:} Ornstein--Uhlenbeck process calibrated to GOOG ---
theta-hat equals 0.0121 ($\hat\theta = 0.0121$),
m-hat equals 0.0207 ($\hat m = 0.0207$),
sigma-hat equals 0.0187 ($\hat\sigma = 0.0187$).
Notice how the paths oscillate around the dashed long-run mean
without diverging to infinity.
The reversion is slow and gentle --- consistent with the half-life of 57 days.

% ================================================================
\section{Slide 13. Empirical GOOG vs Simulated OU Benchmark}
% ================================================================

This is a direct visual comparison.

The orange line is the actual detrended GOOG log-price
($\widetilde{X}_{t_i}^k$) from 2023 to 2026.
The navy line is one Ornstein--Uhlenbeck simulation
with the parameters estimated from GOOG.

The paths are not identical --- they are two different realizations of processes
with similar statistical properties.
The comparison is illustrative and educational, not a statistical test.
But the visual similarity supports the Ornstein--Uhlenbeck interpretation.

% ================================================================
\section{Slide 14. Model Assessment Criterion}
% ================================================================

This slide directly addresses Professor Fatkullin's methodological comment.

The question was: what does ``compatible'' mean?
You cannot just draw two pictures that look similar and say the model fits.
You need an explicit criterion.

The approach used here has three parts.

\textbf{First:} fit parameters using one simple statistic ---
the one-step autoregressive model of order one (AR(1)) regression on detrended levels,
giving alpha-hat ($\hat\alpha$) and phi-hat ($\hat\phi$).

\textbf{Second:} check other implications of the fitted model
that were \emph{not used} in the fitting:
\begin{itemize}
  \item Validation one: does the theoretical decay ($\hat\phi^\ell$)
        match the empirical autocorrelation function at lags 2 through 20?
  \item Validation two: does the model's conditional mean ---
        ($\hat m + \hat\phi^h(\widetilde{X}_t - \hat m)$) ---
        predict the empirical conditional mean at h equals 5 days?
\end{itemize}

\textbf{Third:} report the visible failure ---
time-varying volatility.

This is not a full transition-kernel test, as the professor also suggested.
But it is a systematic diagnostic that gives explicit meaning
to the word ``compatible.''

% ================================================================
\section{Slide 15. Model Assessment: Quantitative Table}
% ================================================================

The table shows all metrics together.

\textbf{Raw R-squared} ($R^2$) is high for all three stocks:
0.883 for NVDA, 0.822 for META, 0.866 for GOOG.
The geometric Brownian motion benchmark applies to raw log-price levels.

\textbf{Ornstein--Uhlenbeck autoregressive R-squared} is approximately 0.98
for all three stocks: the one-step fit is very tight.

\textbf{Autocorrelation function root mean squared error} (ACF RMSE)
is small, especially for GOOG at 0.007:
the theoretical decay curve ($\hat\phi^\ell$) closely tracks
the empirical autocorrelation function.
For NVDA the deviation is larger at 0.049 ---
the series is less cleanly Ornstein--Uhlenbeck-like.

\textbf{Transition root mean squared error} (RMSE)
is approximately 0.005 to 0.007:
the conditional-mean prediction is accurate.

\textbf{Rolling volatility ratio} is greater than 3 to 7 times for all stocks:
this is where the constant-sigma stochastic differential equation clearly fails.
META shows the worst ratio at 6.8 times ---
it had a particularly turbulent period during 2023.

% ================================================================
\section{Slide 16. Main Findings}
% ================================================================

The answer is qualified, not absolute.

\textbf{Geometric Brownian motion (GBM)} is appropriate for raw log-prices ---
these are trend-dominated, growing over time,
and the multiplicative-noise picture matches well.

\textbf{Ornstein--Uhlenbeck (OU)} is appropriate for detrended log-prices ---
these show persistent, slowly decaying autocorrelation
with a half-life of about three months,
and both autocorrelation-decay and conditional-mean validation pass.

\textbf{Constant-sigma diffusion} is useful but limited ---
the Itô language describes daily increments naturally,
but volatility is clearly time-varying.

No single low-dimensional stochastic differential equation
captures all features simultaneously.

% ================================================================
\section{Slide 17. Limitations and Natural Extensions}
% ================================================================

The main limitations are as follows.

Linear detrending may not be the right way to remove the trend;
factor-based detrending could give different results.

The transition-mean check is a simplified diagnostic,
not the full transition-kernel analysis that Professor Fatkullin mentioned.

Constant-volatility stochastic differential equation
misses volatility clustering entirely.

The most natural next step mathematically is
\textbf{realized quadratic variation}:
for a diffusion with drift u and diffusion coefficient v,
the quadratic variation ($\langle X\rangle_t$) ---
converges to ($\int_0^t v^2\,ds$).
In practice, summing squared daily returns gives a direct empirical estimate
of the integrated variance.
This would be a more rigorous diffusion diagnostic than rolling volatility.

Another natural direction is the \textbf{Fokker--Planck equation}
for the Ornstein--Uhlenbeck process,
which gives the evolution of the probability density ---
that would allow a more formal transition-kernel analysis
as suggested by the professor.

% ================================================================
\section{Slide 18. Conclusion}
% ================================================================

The conclusion is deliberately qualified, not absolute.

The research question was: can simple stochastic differential equations
describe equity-market time series?
The answer is: yes, to a meaningful extent,
but only approximately and only for the right transformations of the data.

The methodological lesson is that ``compatible'' is not a visual judgment.
You fit by some statistics and validate by others.

This is an educational mini-research, and I think the value is precisely in making
the connection between Øksendal's theoretical framework and real market data explicit ---
showing which formulas apply where, and what their limitations are.

\end{document}
